\begin{proposition}[Refuted: universal membership of the column module in the proposed resolving class]
\label{A:P1-torsion}
For a small category $A$, let $\mathcal R_A$ be the smallest class of
covariant integral modules containing every $\mathbb Z[H]$, for
$H:A\to\mathrm{Set}$, and closed under extensions, kernels of
epimorphisms between members, arbitrary direct sums, direct summands,
filtered colimits, and $F\mapsto F\otimes_{\mathbb Z}\mathbb Z[H]$.
Every member of $\mathcal R_A$ is objectwise torsion-free.
There is a finite filtered category $A$ and three parallel arrows
$a_0,a_1,a_2:p\to q$ for which
\[
 Q_a=\operatorname{coker}\bigl(P_q\longrightarrow P_p^2,
             \ t\longmapsto(t(a_1-a_0),t(a_2-a_0))\bigr)
 \notin\mathcal R_A.
\]
Thus the first proposed target of approach P1 cannot work uniformly,
even for filtered categories.  This does not refute P1's alternative
of finding a different module with the same $T_1$, and it is not a
separator between totality and exactness.
\end{proposition}
\begin{proof}
Permutation modules are objectwise free.  Subgroups, extensions,
direct sums and direct summands of torsion-free abelian groups are
torsion-free.  A filtered colimit is torsion-free because an equation
$n[x]=0$, $n\ne0$, holds at some later stage, where it implies that
the representative itself is zero.  Tensoring with a free abelian
group is a direct sum of copies.  All operations are evaluated
objectwise, so the class of objectwise torsion-free modules contains
$\mathcal R_A$.

Here is the promised category completely.  Its objects are $p,q,x,t$.
Besides identities, its arrows are
\[
 a_0,a_1,a_2:p\to q,\qquad
 g_0,g_1,g_2,g_3:q\to x,\qquad
 c_0,c_1,c_2:p\to x,
\]
and one arrow $u_y:y\to t$ for each $y=p,q,x$.  There are no other
arrows.  Identities compose as identities, every composition ending
at $t$ is the indicated unique arrow, and the other compositions are
specified by
\[
\begin{array}{c|ccc}
 &a_0&a_1&a_2\\\hline
g_0&c_0&c_1&c_2\\
g_1&c_0&c_2&c_1\\
g_2&c_1&c_0&c_2\\
g_3&c_1&c_2&c_0
\end{array}
\]
Associativity holds: a composable triple of nonidentity arrows must
follow $p\to q\to x\to t$, and both composites are $u_p$.
All other associativity equations contain an identity.
The object $t$ is terminal, so the category is nonempty, connected
and filtered.  It is total directly: for any $H:A\to\mathrm{Set}$,
the object $(a,h)$ of its category of elements maps to
$(t,H(u_a)h)$.  Its connected component is exactly the fibre of this
value in $H(t)$ and has that terminal object; the assertion is also
valid if $H(t)$ is empty.

At $x$, $P_q(x)=\mathbb Z^4$ on the $g_j$ and $P_p(x)=\mathbb Z^3$
on the $c_i$.  The column map has its image in
$\ker(\operatorname{aug}:\mathbb Z^3\to\mathbb Z)^2$, a rank-four
direct summand of $\mathbb Z^6$.  With bases $c_1-c_0,c_2-c_0$ in
both augmentation kernels, its matrix is
\[
 D=\begin{pmatrix}
 1&0&-1&-1\\
 0&1&0&1\\
 0&1&-1&-1\\
 1&0&1&0
 \end{pmatrix}.
\]
Its determinant is $3$, and its upper-left $3\times3$ minor has
determinant $-1$.  Hence its Smith invariant factors are
$1,1,1,3$, and
\[
 Q_a(x)\cong\mathbb Z^2\oplus\mathbb Z/3.
\]
More explicitly, $v=(1,0,1,0)^{\mathrm t}$ represents a nonzero
element of order $3$ in the reduced cokernel: one has
$D(1,1,-1,-1)^{\mathrm t}=3v$, while the unique rational preimage
of $v$ is $(1,1,-1,-1)^{\mathrm t}/3$, which is not integral.
Thus $Q_a$ cannot belong to $\mathcal R_A$.
The program \texttt{code/algebra\_column\_torsion.py} reconstructs
the category, checks all $52$ composable pairs and $121$ composable
triples, verifies the identities and terminal object, and recomputes
the Smith form and this order-three certificate.  Its output is
\texttt{code/algebra\_column\_torsion.out}.  The original colimit is
evaluation at $t$, so $T_1(Q_a)=0$ in this example; no nonexactness is
being inferred from its objectwise torsion.
\end{proof}
